% Compile twice: lualatex 20001258.tex
\documentclass[11pt]{article}
\usepackage[margin=25mm]{geometry}
\usepackage{fontspec,unicode-math}
\setmainfont{Latin Modern Roman}
\setmathfont{Latin Modern Math}
\usepackage{amsmath,amssymb,mathtools}
\usepackage{xurl,hyperref}
\hypersetup{colorlinks=true,urlcolor=blue,linkcolor=blue}
\setcounter{secnumdepth}{0}
\setlength{\parindent}{0pt}
\setlength{\parskip}{5pt}
\setlength{\emergencystretch}{3em}
\newcommand{\sref}[2]{\hyperlink{source:#1:#2}{[S#2]}}
\newcommand{\cataloguescope}[1]{\paragraph{Source question.}#1}
\newcommand{\floor}[1]{\left\lfloor#1\right\rfloor}
\begin{document}
% UnsolvedMath ID 20001258; source catalogue code AIM-COMPUTATION-0096
\setcounter{section}{20001257}
\section{Information and conditioning audit for zeta-zero primality and factoring proposals}\label{20001258}
{\small\sffamily UnsolvedMath ID 20001258 · Computer Science}
\subsection{Definitions and mathematical statement}

\paragraph{Shared notation.}Standard number systems are denoted by $\mathbb N,\mathbb Z,\mathbb Q,\mathbb R,\mathbb C$; index ranges, computational models and specialized notation follow the cited problem source and the definitions below.


The Riemann zeta function $\zeta(s)$ has nontrivial complex zeros $\rho$. The von Mangoldt function satisfies $\Lambda(n)=\log p$ for a prime power $n=p^k$ and is zero otherwise. The source suggests extracting primality or factorization information from sums over zeros with imaginary parts in $[T,2T]$.
\paragraph{Statement recorded in the supplied source.}
Question 8. Look at ∑ nρ over the zeros ρ of the Riemann zeta function with the imaginary part of ρ the interval [ T, 2T ]. This is ≈ T / (2 π)Λ( n). Can you make a primality test out of this? (Conrey) The sum ∑ nρζ′(ρ)/ζ ′′ (ρ) ≈ T / (2 π) log p log q if n = pq. Can you make a factoring algorithm out of this?

Problem/ \sref{20001258}{2}
\paragraph{Scope qualification.}
The source's displayed approximate identities are motivations, not asserted error bounds. Computing zeros accurately, treatment of multiplicities, truncation, and bit complexity remain unspecified.
\paragraph{Source status.}
The supplied catalogue labels this record partially solved. Its source-status metadata was checked against the supplied snapshot on 2026-10-11. This import does not independently establish current open status or mathematical correctness.
\subsection{Short English statement}
Can numerical information about zeta zeros be turned into a primality test or factoring algorithm?
\subsection{Sources}
\begin{itemize}
\item[\hypertarget{source:20001258:1}{S1}] Ulam AI and the UnsolvedMath Contributors, UnsolvedMath: A Curated Collection of Open Mathematics Problems, record 20001258 (AIM-COMPUTATION-0096). Catalogue curation and metadata: CC BY 4.0. The supplied source question is retained; definitions and the English summary are editorial additions.
\url{https://huggingface.co/datasets/ulamai/UnsolvedMath}.
\item[\hypertarget{source:20001258:2}{S2}] AIM Workshop Problem Lists, original problem source linked by catalogue record AIM-COMPUTATION-0096
\url{https://aimath.org/WWN/primesinp/primesinp.pdf}.
\end{itemize}
\paragraph{Preparation and review.}
Original question attributed to the cited source and UnsolvedMath contributors. Codex (GPT-6; exact serving revision not exposed) assisted with definitions, English summaries, source inspection and formatting. No new human mathematical authorship or independent proof verification is asserted. Canonical ID 20001258 and source code AIM-COMPUTATION-0096 checked against the complete supplied registry and existing repository statements; known flagged repetitions and obvious identical conjectures excluded. Full mathematical-equivalence/literature review has not been performed. Statement preparation and final review followed contribution prompts C and D; automated checks verify structure and identity, not open status.
\label{end:20001258}
\end{document}
